\begin{itemize}
    \item \href{https://github.com/pankajgarkoti/magi}{mAGI} \\
    mAGI is an agent-first desktop built on Fedora and Hyprland: one orchestrator agent in tmux runs the desktop with worker agents, an event bus, and a shared browser.
    \item \href{https://computer-reinvention.github.io}{trie} \\
    trie is a symbol-level codebase indexing layer that exposes an MCP server interface (\texttt{trie\_grep}, \texttt{trie\_read}, \texttt{trie\_trace}), enabling AI coding agents to navigate a codebase by intent rather than brute-force file search. In benchmarks on the same question, agents with trie used 3 times less tool calls. The index bootstraps itself in the project and is incrementally kept in sync with the source code. In active development.
    \item \href{https://github.com/pankajgarkoti/daily-notes.nvim}{daily-notes.nvim} \\
    daily-notes.nvim is a lightweight Neovim plugin that replicates Obsidian's markdown template based daily note functionality inside Neovim. The plugin by default carries over the contents of previous day's note, headers that should not be copied can be specified separately. The plugin is highly configurable and allows custom datetime formats and paths etc.
    \item \href{https://github.com/pankajgarkoti/chainfactory}{ChainFactory} \\
    ChainFactory is a declarative system for building multi-step LLM workflows in a YAML-like syntax. Chains are composed of links whose outputs feed into subsequent inputs, with sequential and parallel execution, auto-generated prompts from purpose statements, and type-safe structured output via Pydantic models. Ideal for split-map-reduce style inference pipelines.
    \item Computer Reinvention \\
    My personal open-source lab exploring how LLMs can reshape the way developers read, navigate, and reason about code. All work is public and built in the open.
\end{itemize}
